\newpage
\section{CONCLUSION}

In this project, a prototype of an automatic directional antenna alignment solution utilizing reinforcement learning (RL) and low-cost embedded hardware is described. The proposed design combines a specially developed mechanical framework with sensor and actuator control using ESP32-S3, Q-Learning without online learning, simulation analysis, and testing of a fully functional model into one end-to-end solution. From the final experiments, the conclusion may be drawn that for all analyzed cases, the solution based on reinforcement learning provides satisfactory convergence and good quality of alignment at the same time. In contrast to exhaustive scanning, the RL controller converges significantly faster and achieves similar RSSI values compared to the reference point. The proposed design works better than hill climbing in real-world situations.

The design of the system was influenced by costs and availability considerations for particular components, such as the custom PCB slip ring, which was included in the design, and the selection of actuators that were available on the market. This explains why the current prototype demonstrates deficiencies in terms of durability, signal quality, accuracy, and the range of test conditions.

Moreover, the RL controller is trained offline in a discrete state space environment, and its operation might be compromised when there is a significant difference between the operational conditions and the training set.

In conclusion, the project has demonstrated that reinforcement learning can be employed as an effective method for autonomous antenna pointing in affordable wireless devices. This paper provides initial groundwork for future investigation into beam steering, intelligent radio channels, and embedded RL control algorithms.